%%
%% fall-lecture02.tex
%% 
%% Made by Alex Nelson <pqnelson@gmail.com>
%% Login   <alex@lisp>
%% 
%% Started on  2026-10-01T10:37:16-0700
%% Last update 2026-10-01T10:37:16-0700
%% 

\lecture{}

\begin{node}
Recall, we looked at 3 conditions for using the ocntinuum
approximation:
\begin{enumerate}
\item dilute gas $\delta\gg d$
\item physical length scale is small $\Kn=\lambda/L\ll1$
\item meaningful statistics $nV\gg1$.
\end{enumerate}
We will finish discussing meaningful statistics today.
\end{node}

\begin{node}[Roadmap for today's lecture]
  There are also several things we want discuss:
  \begin{enumerate}
  \item Maxwell--Boltzmann distribution
  \item Molecular fluxes across surfaces.
  \end{enumerate}
\end{node}

\begin{node}[Meaningful statistics]
If you are measuring a quantity but its standard deviation is large
enough, then you need to be careful. The example given was the mass
density. We average over some small-but-finite volume $V$ and examine
the probability density $P(N)$ for the probability of $N$ particles to
be in the volume $V$. This may be approximated using the Poisson
distribution with parameter $nV$ where $n$ is the number density
(number of particles per unit volume). If we take the limit
$nV\to\infty$ by taking $V\to\infty$, then the probability
distribution becomes the $\delta$ distribution (i.e., sharply peaks)
\begin{equation}
P(N/V)\to\delta(N/V).
\end{equation}
You want the standard deviation to be small relative to the mean. The
intuition justifying this is: if we sample around the neighborhood,
then we should expect small fluctiations. We could trake
\begin{equation}
\frac{L}{\delta}=V^{1/3}\frac{1}{(1/n)^{1/3}}=(nV)^{1/3}\geq\mbox{threshold},
\end{equation}
where the threshold is usually taken to be 100 or so, giving us the heuristic
\begin{equation}
nV\geq10^{6}.
\end{equation}
\end{node}

\begin{node}
So, to recap, we have three inequalities we want to be satisfied:
\begin{enumerate}
\item $\delta/d\gg1$ for the dilute gas limit
\item $\Kn=\lambda/L\ll1$ for the physical length scale to be large enough
\item $L/\delta\gg1$ for meaningful statistics
\end{enumerate}
Traditionally, we have a plot of physical length scale $L$ against
$\delta/d$, and we plot three inequalities (depending on our heuristic
used)
\begin{enumerate}
\item $\delta/d\geq7$ for the dilute gas limit
\item $\Kn=\lambda/L\leq0.1$ for the physical length scale to be large enough
\item $L/\delta\geq10^{2}$ for meaningful statistics.
\end{enumerate}
Recall Eq~\eqref{eq:fall:lec1:mean-free-path}, the mean free path for
a Maxwell distribution has $\lambda\sim\delta^{3}/d^{2}$, so
$\lambda/L\sim(\delta/d)^{3}(d/L)$. We also see $L/\delta=(L/d)(d/\delta)$.
This gives us a way to express the inequalities using factors of $L/d$
and $\delta/d$ alone. Traditionally, we plot these inequalities
using logarithmic scales (see, e.g., Ch~2 of \textit{Heat Transfer and Fluid Flow in Minichannels and Microchannels},
\doi{10.1016/C2011-0-07521-X}; or G.A.~Bird~\cite[Fig~1.6]{bird1976molecular}
or his later book \textit{Molecular Gas Dynamics and the Direct Simulation of Gas Flows}, Oxford University
press (1998)). This means that we have the inequalities
\begin{enumerate}
\item $\log(\delta/d)\geq\log(7)$ for dilute gases (and $\leq\log(7)$
for dense gases)
\item $\log(\Kn)=3\log(\delta/d)-\log(L/d)\leq\log(0.1)$ for the
Navier-Stokes equations being valid (and $\geq\log(0.1)$ for the
microscopic approach being necessary)
\item $\log(L/\delta)=\log(L/d)-\log(\delta/d)\geq2\log(10)$ for
statistical homogeneity, i.e., ``insignificant statistical fluctuations''
(and $\leq2\log(10)$ for significant statistical fluctuations where
inhomogeneities are noticeable). 
\end{enumerate}
When $y=\log(L/d)$ and $x=\log(\delta/d)$ and we work with the common
base-10 logarithm, these are the system of inequalities
\begin{subequations}
  \begin{align}
    x&\geq0.845\\
    3x-y&\leq-1\\
    y-x&\geq2,
  \end{align}
\end{subequations}
or
\begin{subequations}
  \begin{align}
    x&\geq0.845\\
    y&\geq 3x+1\\
    y&\geq x+2.
  \end{align}
\end{subequations}
\end{node}

\subsection{Microscopic approach}

\begin{node}
The miscroscopic approach treats fluids as a collection of
particles. There are different ways to this approach:
\begin{enumerate}
\item Molecular dynamics: treat each molecule as a rigid sphere,
  obeying Newtonian mechanics;
\item Direct sumlation Monte Carlo methods (DSMC) looks at probability
  distributions to solve the Boltzmann equations;
\item Lattice Boltzmann method --- a ``mesoscale'' method where we
  calculate interactions beween groups of molocules.
\end{enumerate}
\end{node}

\begin{node}
Consider a dilute gas at equilibrium (which is when the ``meaningful
statistics'' condition is satisfied). Each particle has a molecular
velocity $\vec{c}$ and a mean velocity
\begin{equation}
\vec{u}=(\langle c_{1}\rangle, \langle c_{2}\rangle, \langle c_{3}\rangle).
\end{equation}
We can look at the distribution $f_{0}(\vec{c})$ of velocity about the
mean velocity. This is the \define{Maxwell distribution}
\begin{equation}
f_{0}(\vec{c})=\left(\frac{m}{2\pi\Boltzmann T}\right)^{3/2}\exp\left[\frac{-m}{2\Boltzmann T}|\vec{u}-\vec{c}|^{2}\right]
\end{equation}
where $\Boltzmann$ is the Boltzmann constant (which cancels all units in
the exponential), $m$ is the molar mass, and $T$ is the temperature.
Observe that the Maxwell distribution is ``just'' the Gaussian
distribution.

A lot of microscopic approaches use the Maxwell distribution to
initialize or warm-up.
\end{node}

\begin{node}
What are some properties of the Maxwell distribution?
\begin{enumerate}
\item $\iiint_{\RR^{3}}f_{0}(\vec{c})\,\D^{3}\vec{c}=1$ it's a
  probability distribution
\item We assume that it can factorize into a product of identical
  distributions $f_{0}(\vec{c})=P(c_{1})P(c_{2})P(c_{3})$
\item $u_{1}=\langle c_{1}\rangle=\iiint_{\RR^{3}}c_{1}f_{0}(\vec{c})\,\D^{3}\vec{c}=\int_{\RR}c_{1}P(c_{1})\,\D{c_{1}}$
  is the mean, and
  \begin{subequations}
    \begin{align}
\langle(c_{1}-u_{1})^{2}\rangle
&=\langle(c'_{1})^{2}\rangle\\
&=\iiint_{\RR^{3}}(c_{1}-u_{1})^{2}f_{0}(\vec{c})\,\D^{3}\vec{c}\\
&=\frac{\Boltzmann T}{m}
    \end{align}
  \end{subequations}
  is the variance;
\item Equipartition of energy: $\langle(c'_{1})^{2}\rangle=\langle(c'_{2})^{2}\rangle=\langle(c'_{3})^{2}\rangle$
  variance of each component of velocity is the same (by independence
  of components). Think of the kinetic energy relative to the mean in
  each dimension
  \begin{equation}
KE_{j}=\frac{1}{2}m\langle(c'_{j})^{2}\rangle
  \end{equation}
  so
  \begin{equation}
KE = KE_{1}+KE_{2}+KE_{3}=\frac{3}{2}\Boltzmann T.
  \end{equation}
\item Equipartition of energy for rotational and vibrational
  energies. This gives us
  \begin{equation}
E = (\mbox{degrees of freedom})\cdot\frac{1}{2}\Boltzmann T.
  \end{equation}
\end{enumerate}
\end{node}

\begin{node}
If we look at the expected value of the speed, we find (changing to
spherical coordinates where the radial component is precisely $|\vec{c}|$):
\begin{subequations}
  \begin{align}
\langle|\vec{c}|\rangle
&=\iiint_{\RR^{3}}|\vec{c}|f_{0}(\vec{c})\,\D^{3}\vec{c}\\
&=\int^{\pi}_{0}\int^{2\pi}_{0}\int^{\infty}_{0}cf_{0}(\vec{c})\cdot
c^{2}\sin(\theta)\,\D c\,\D\phi\,\D\theta\\
&=\sqrt{\frac{8\Boltzmann T}{\pi m}}.
  \end{align}
\end{subequations}
\end{node}

\begin{remark}
For real gases, a better approximation would be to make the
replacement
\begin{equation}
\frac{\Boltzmann T}{m}\mapsto\frac{p V_{m}}{M}
\end{equation}
where $p$ is the pressure, $V_{m}$ is the molar volume, and $M$ is the
molar mass. See~\cite{lente2025direction} for a Mathematical
derivation. But there does not seem to be any experiments done to
empirically test the goodness-of-fit of this distribution.
\end{remark}

\subsection{Molecular fluxes across a surface}

\begin{node}
Consider the number of molecules crossing a region $\D S$ in a time interval
$\D t$ for a given velocity $\vec{c}$. We can look at the tube of height
$\vec{c}\cdot\UnitNormalVec{n}\,\D t$ and the flux is
\begin{subequations}
  \begin{align}
    \D\widehat{N} &= \begin{pmatrix}\mbox{number of}\\
      \mbox{molecules per}\\
      \mbox{unit volume}
    \end{pmatrix}\begin{pmatrix}\mbox{volume of}\\
      \mbox{cylinder}
    \end{pmatrix}\\
    &=(nf_{0}(\vec{c})\,\D^{n}c)\cdot(\D S\,\vec{c}\cdot\UnitNormalVec{n}\,\D t)
  \end{align}
\end{subequations}
which counts the number of molecules with velocity $\vec{c}$ passing
through the cylinder. When each molecule carries some physical
quantity $Q$, we can then consider the $Q$-flux. For some examples:
\begin{enumerate}
\item $Q=1$ for the number of molecules cross $\D S$
\item $Q=\vec{c}$ for the velocity of molecules crossing $\D S$
\item $Q=\frac{1}{2}m\vec{c}\cdot\vec{c}$ for the kinetic energy of
  molecules crossing $\D S$.
\end{enumerate}
Then
\begin{equation}
\D Q=Q(\vec{c})\,\D\widehat{N},
\end{equation}
giving us
\begin{equation}
\mbox{$Q$-flux}=\iiint_{\RR^{3}}\D Q=\iiint_{\RR^{3}}nQ(\vec{c})f_{0}(\vec{c})\vec{c}\cdot\UnitNormalVec{n}\,\D^{3}c,
\end{equation}
which describes the amount of $Q$-flux moved per unit time per unit
area. This describes \emph{both} incoming \emph{and} outgoing molecules.
\end{node}

% This was at the start of the next lecture.
\begin{node}[One way flux]
We can write 
\begin{equation}
\mbox{$Q$-flux}=\langle nQ(\vec{c})\vec{c}\cdot\UnitNormalVec{n}\rangle
\end{equation}
for the total $Q$-flux.

We can also look at the ``one-way flux'' when $\vec{c}\cdot\UnitNormalVec{n}>0$
describing the incoming flux
\begin{equation}
\langle nQ(\vec{c})\vec{c}\cdot\UnitNormalVec{n}\rangle^{+}
=\langle\delta(\sgn(\vec{c}\cdot\UnitNormalVec{n})=1)nQ(\vec{c})\vec{c}\cdot\UnitNormalVec{n}\rangle
\end{equation}
The outgoing flux is similarly when $\vec{c}\cdot\UnitNormalVec{n}<0$,
\begin{equation}
\langle nQ(\vec{c})\vec{c}\cdot\UnitNormalVec{n}\rangle^{-}
=\langle\delta(\sgn(\vec{c}\cdot\UnitNormalVec{n})=-1)nQ(\vec{c})\vec{c}\cdot\UnitNormalVec{n}\rangle
\end{equation}
These can be combined to give us the total $Q$-flux:
\begin{equation}
\langle nQ(\vec{c})\vec{c}\cdot\UnitNormalVec{n}\rangle
=\langle nQ(\vec{c})\vec{c}\cdot\UnitNormalVec{n}\rangle^{+}
+\langle nQ(\vec{c})\vec{c}\cdot\UnitNormalVec{n}\rangle^{-}.
\end{equation}
\end{node}

\begin{remark}
We are assuming we have a dilute gas in equilibrium. Here, ``in equilibrium''
means the distribution $f_{0}(\vec{c})$ does not change with respect
to time.
\end{remark}

\begin{remark}
The reference for the first and second lecture seems to be based
heavily on Bird~\cite[Ch~1~\&~3]{bird1976molecular}.
\end{remark}
